\documentclass[11pt]{article}
\usepackage[utf8]{inputenc}
\usepackage[T1]{fontenc}
\usepackage{amsmath,amssymb,mathtools}
\usepackage[margin=1in]{geometry}
\usepackage{hyperref}
\usepackage{array}
\usepackage{parskip}
\setlength{\parindent}{0pt}
\hypersetup{colorlinks=true,linkcolor=blue,urlcolor=blue}
\title{Lecture 11: Frequency Response and Filtering}
\author{}
\date{}
\begin{document}
\maketitle

\section*{Frequency Response}

An LTI system acts as a \textbf{frequency-by-frequency multiplier}. For a periodic input with FS coefficients $\{a_k\}$, the output coefficients are:

\[b_k = a_k \cdot H(jk\omega_0)\]

where $H(j\omega)$ is the \textbf{frequency response} (eigenvalue at frequency $\omega$).

\subsection*{CT Frequency Response}

\[H(j\omega) = \int_{-\infty}^{\infty} h(\tau) e^{-j\omega\tau} d\tau\]

For causal $h(t) = e^{-at}u(t)$:

\[H(j\omega) = \frac{1}{a + j\omega}\]

\subsection*{DT Frequency Response}

\[H(e^{j\omega}) = \sum_{n=-\infty}^{\infty} h[n] \, e^{-j\omega n}\]

\section*{Output of LTI System to Periodic Input}

Given input $x(t) = \sum a_k e^{jk\omega_0 t}$:

\[y(t) = \sum_{k=-\infty}^{\infty} a_k \, H(jk\omega_0) \, e^{jk\omega_0 t}\]

Each harmonic is independently scaled by $|H(jk\omega_0)|$ and phase-shifted by $\angle H(jk\omega_0)$.

\section*{Filter Types}

\begin{center}
\begin{tabular}{|l|l|l|}
\hline
\textbf{Filter} & \textbf{Passes} & \textbf{Attenuates} \\
\hline
\textbf{Lowpass} & Low frequencies (near $\omega = 0$) & High frequencies \\
\hline
\textbf{Highpass} & High frequencies & Low frequencies (near $\omega = 0$) \\
\hline
\textbf{Bandpass} & Frequencies in a specific band & Frequencies outside the band \\
\hline
\end{tabular}
\end{center}

\section*{Ideal vs Real Filters}

\textbf{Ideal filters:} Perfectly rectangular magnitude response (1 in passband, 0 in stopband). Physically \textbf{unrealizable} (noncausal).

\textbf{Real filters:} Gradual roll-off with a \textbf{transition band} between passband and stopband.

\section*{RC Circuit Example}

\subsection*{Lowpass: $H(j\omega) = \frac{1}{1 + j\omega RC}$}

\begin{itemize}
  \item Cutoff frequency: $\omega_c = 1/RC$
  \item At $\omega_c$: magnitude drops to $1/\sqrt{2}$ of DC value (-3 dB point)
  \item Phase at cutoff: $-45^\circ$
\end{itemize}

\subsection*{Highpass: $H(j\omega) = \frac{j\omega RC}{1 + j\omega RC}$}

\begin{itemize}
  \item Same cutoff frequency $\omega_c = 1/RC$
  \item At DC ($\omega = 0$): output is zero
  \item At high frequencies: output approaches input
\end{itemize}

\section*{General First-Order Lowpass: $H(j\omega) = \frac{1}{a + j\omega}$}

\begin{itemize}
  \item Cutoff: $\omega_c = a$ (the pole location)
  \item $|H(j\omega_c)| = \frac{1}{\sqrt{2}} |H(0)|$ (definition of -3 dB)
  \item Increasing $a$ widens bandwidth (less aggressive filtering)
\end{itemize}

\section*{FIR vs IIR Filters (DT)}

\begin{center}
\begin{tabular}{|l|l|l|l|}
\hline
\textbf{Type} & \textbf{Impulse Response} & \textbf{Memory} & \textbf{Example} \\
\hline
\textbf{FIR} (Finite Impulse Response) & Finite duration & Finite & Moving average \\
\hline
\textbf{IIR} (Infinite Impulse Response) & Infinite duration & Infinite & Recursive filter \\
\hline
\end{tabular}
\end{center}

\section*{Time-Frequency Trade-off}

Narrower filters (sharper frequency selection) require longer impulse responses (more time-domain spread). You cannot have both sharp frequency cutoff and compact time response:

\[\omega_c \cdot t_r \approx \text{constant}\]

\end{document}
